% PRISMA-ScR 22-item checklist
% Adapted from: Tricco et al. (2018) Ann Intern Med 169(7):467–473
% Fill in "Reported on page" column as each section is drafted.
\begingroup
\renewcommand{\arraystretch}{1.2}
\footnotesize
\begin{longtable}{p{0.4cm} p{1.6cm} p{8.5cm} p{1.2cm}}
\caption{PRISMA-ScR 22-item checklist}\\
\toprule
\textbf{\#} & \textbf{Section} & \textbf{Item} & \textbf{Page} \\
\midrule
\endfirsthead
\multicolumn{4}{c}{\tablename\ \thetable{} --- continued}\\
\toprule
\textbf{\#} & \textbf{Section} & \textbf{Item} & \textbf{Page} \\
\midrule
\endhead
\midrule
\multicolumn{4}{r}{Continued on next page}\\
\endfoot
\bottomrule
\endlastfoot
% --- Title and Abstract ---
1 & Title &
  Identify the report as a scoping review. &
  \annote{Title} \\
2 & Abstract &
  Provide a structured summary including background, objectives,
  eligibility criteria, sources of evidence, charting methods,
  results, and conclusions. &
  \annote{Abstract} \\
% --- Introduction ---
3 & Introduction &
  Describe the rationale for the review in the context of what
  is already known. &
  \annote{Sec. 1} \\
4 & Introduction &
  Provide an explicit statement of the questions and objectives
  being addressed with reference to key concepts. &
  \annote{Sec. 1 (RQ)} \\
% --- Methods ---
5 & Methods &
  Indicate whether a review protocol exists; describe provisions
  for protocol amendments. &
  \annote{Sec. 3.1} \\
6 & Methods &
  Specify characteristics of sources of evidence used as
  eligibility criteria (e.g., years considered, language,
  publication status). &
  \annote{Sec. 3.2} \\
7 & Methods &
  Describe all information sources (e.g., databases with dates
  of coverage, contact with authors). &
  \annote{Sec. 3.3} \\
8 & Methods &
  Present the full electronic search strategy for at least one
  database, including any limits used. &
  \annote{Sec. 3.3 + App. B} \\
9 & Methods &
  Describe the process for selecting sources of evidence. &
  \annote{Sec. 3.4} \\
10 & Methods &
  Describe the data charting process (e.g., piloted forms,
  independent reviewers). &
  \annote{Sec. 3.5} \\
11 & Methods &
  List and define all variables for which data were sought. &
  \annote{Sec. 3.5} \\
12 & Methods &
  If critical appraisal of sources was performed, describe the
  rationale and methods used. &
  \annote{Sec. 3.6} \\
13 & Methods &
  Describe the methods of charting data and synthesising
  results. &
  \annote{Sec. 3.7} \\
% --- Results ---
14 & Results &
  Give numbers of sources screened, assessed for eligibility,
  and included in the review, with reasons for exclusions at
  each stage, ideally using a flow diagram. &
  \annote{Sec. 4.1 + Fig. 1} \\
15 & Results &
  For each source of evidence, present characteristics for
  which data were charted. &
  \annote{Sec. 4.1} \\
16 & Results &
  If critical appraisal was performed, present data and
  describe implications. &
  \annote{Sec. 4.1} \\
17 & Results &
  For each included source, present the data that were charted. &
  \annote{Sec. 4.1} \\
18 & Results &
  Summarise and present the charting results with reference to
  the review questions and objectives. &
  \annote{Sec. 4.2--4.4} \\
% --- Discussion ---
19 & Discussion &
  Summarise the main results, including key concepts and
  evidence gaps. &
  \annote{Sec. 5} \\
20 & Discussion &
  Discuss limitations of the scoping review process. &
  \annote{Sec. 5.1} \\
21 & Discussion &
  Provide a general interpretation of the results with respect
  to the review questions and objectives. &
  \annote{Sec. 5} \\
% --- Other ---
22 & Other &
  Describe sources of funding for the review. &
  \annote{Funding stmt} \\
\end{longtable}
\endgroup
